%------------------------------------------------------------------------------%
\documentclass[fleqn,utf8,aspectratio=169,12pt]{beamer}

\usepackage{gaal}

\title{Norma de um Vetor}

%------------------------------------------------------------------------------%
\begin{document}

\addtitle

%------------------------------------------------------------------------------%
\section{Norma de um Vetor}

%------------------------------------------------------------------------------%
\begin{frame}{Norma de um Vetor}
  \emph{Norma}, \emph{Comprimento} ou \emph{Módulo} do vetor

  \pause
  é o comprimento do segmento orientado
  \pause
  \[
    \norm{\vec{u}}
    \pause
    = \abs{\vec{u}}
    \pause
    = \abs{\overrightarrow{AB}}
  \]

  \pause
  Direção e sentido não interferem na norma

  \pause
  A norma é uma grandeza escalar (número real) não negativa
  \[
    \pause
    \norm{\vec{u}} \in \R
    \qquad\qquad
    \norm{\vec{u}} \geq 0
  \]
\end{frame}

%------------------------------------------------------------------------------%
\framedgraphic*{Interpretação Geométrica}{E-04-Norma-1}{}
\framedgraphic*{Interpretação Geométrica}{E-04-Norma-2}{}
\framedgraphic*{Interpretação Geométrica}{E-04-Norma-3}{}

%------------------------------------------------------------------------------%
\begin{frame}[t]{Cálculo da Norma}
  \vspace{1\baselineskip}

  Para um vetor no plano $\vec{u} \in \R^2$
  \[
    \pause
    \vec{u}
    \pause
    = ( u_1, u_2 )
    \pause
    = \vetor{u_1}{u_2}
  \]

  \pause
  Pelo Teorema de Pitágoras
  \[
    \pause
    \norm{\vec{u}}
    \pause
    = \sqrt{ u_1^2 + u_2^2 }
  \]
\end{frame}

%------------------------------------------------------------------------------%
\begin{frame}[t]{Cálculo da Norma}
  \vspace{1\baselineskip}

  Para um vetor no espaço $\vec{u} \in \R^3$
  \[
    \pause
    \vec{u}
    = ( u_1, u_2, u_3 )
    \pause
    = \vetortri{u_1}{u_2}{u_3}
  \]

  \pause
  Pelo Teorema de Pitágoras
  \[
    \pause
    \norm{\vec{u}}
    \pause
    = \sqrt{ u_1^2 + u_2^2 + u_3^2 }
  \]
\end{frame}

%------------------------------------------------------------------------------%
\begin{frame}[t]{Cálculo da Norma}
  \vspace{1\baselineskip}

  De forma geral, $\vec{u} \in \R^n$
  \[
    \pause
    \vec{u}
    \pause
    = ( u_1, u_2, \dots, u_n )
    \pause
    = \begin{pmatrix}[c]
      u_1 \\
      u_2 \\
      \vdots \\
      u_n
    \end{pmatrix}
  \]
  \pause
  Pelo Teorema de Pitágoras
  \[
    \pause
    \norm{\vec{u}}
    \pause
    = \sqrt{ u_1^2 + u_2^2 + \cdots + u_n^2 }
    \pause
    = \sqrt{ \; \sum_{i=1}^{n} u_i^2 \; }
  \]
\end{frame}

%------------------------------------------------------------------------------%
\input{ex/E-04-Norma-ex-1}

%------------------------------------------------------------------------------%
\begin{frame}{Propriedades da Norma}
  Para quaisquer vetores $\vec{u}$ e $\vec{v}$ de mesma dimensão
  \pause
  \[
    \norm{\vec{u}} \geq 0
  \]

  \pause
  \[
    \norm{\vec{u}} = 0
    \qquad\Leftrightarrow\qquad
    \vec{u} = \vec{0}
  \]

  \pause
  \[
    \norm{-\vec{u}} = \norm{\vec{u}}
  \]

  \pause
  \[
    \norm{\vec{u}+\vec{v}}
    \leq
    \norm{\vec{u}} + \norm{\vec{v}}
    \pause
    \qquad
    \text{\emph{desigualdade triangular}}
  \]
\end{frame}

%------------------------------------------------------------------------------%
\begin{frame}{Relação Entre Norma e Produto por Escalar}
  O norma do produto por escalar satisfaz
  \[
    \norm{\lambda\vec{u}} = \abs{\lambda} \; \norm{\vec{u}}
  \]
\end{frame}

%------------------------------------------------------------------------------%
\input{ex/E-04-Norma-ex-2}
\input{ex/E-04-Norma-ex-3}

%------------------------------------------------------------------------------%
\section{Vetor Unitário}

%------------------------------------------------------------------------------%
\begin{frame}{Vetor Unitário}
  Um \emph{vetor unitário} tem norma igual a 1

  \pause
  Dado um vetor $\vec{v}$

  \pause
  podemos criar um vetor unitário com mesma direção e sentido

  \pause
  dividindo $\vec{v}$ pela sua norma
  \[
    \hat{v} = \frac{\vec{v}}{\norm{\vec{v}}}
  \]
\end{frame}

%------------------------------------------------------------------------------%
%------------------------------------------------------------------------------%
\begin{question}
  Encontre o vetor unitário na mesma direção e sentido do vetor
  \[
    \vec{u} = \vetortri{1}{3}{-2}
  \]
\end{question}

%------------------------------------------------------------------------------%
\begin{resolution}{Solução}
\begin{columns}[t]
\column{0.4\textwidth}
  \onslide<+->{Norma do vetor $\vec{u}$}
  \begin{align*}
    \onslide<+->{\norm{\vec{u}}}
    \onslide<+->{& = \norm{\vetortri{1}{3}{-2}} \\}
    \onslide<+->{& = \sqrt{ 1^2 + 3^2 + (-2)^2} \\}
    \onslide<+->{& = \sqrt{ 1 + 9 + 4} \\}
    \onslide<+->{& = \sqrt{14}}
  \end{align*}
\column{0.6\textwidth}
  \onslide<+->{Vetor unitário}
  \begin{align*}
    \onslide<+->{\hat{u} }
    \onslide<+->{= \frac{\vec{u}}{\norm{\vec{u}}} }
    \onslide<+->{= \frac{1}{\sqrt{14}}\vetortri{1}{3}{-2} }
    \onslide<+->{= \Vetortri{\frac{1}{\sqrt{14}}}{\frac{3}{\sqrt{14}}}{\frac{-2}{\sqrt{14}}}}
  \end{align*}
\end{columns}
\end{resolution}

%------------------------------------------------------------------------------%
\endinput



\input{ex/E-04-Norma-ex-5}
\input{ex/E-04-Norma-ex-6}

%------------------------------------------------------------------------------%
%\section{Lista Mínima}
%
%%------------------------------------------------------------------------------%
%\begin{frame}{Lista Mínima}
%  \begin{enumerate}
%    \item
%  \end{enumerate}
%
%  \vfill
%  Atenção: A prova é baseada no livro, não nas apresentações
%\end{frame}

%------------------------------------------------------------------------------%
\end{document}
%------------------------------------------------------------------------------%


